% ===================== PREAMBULE COMMUN (à copier VERBATIM en tête de CHAQUE .tex) =====================
\documentclass[11pt,a4paper]{article}
\usepackage[utf8]{inputenc}
\usepackage[T1]{fontenc}
\usepackage{lmodern}
\usepackage[french]{babel}
\usepackage{amsmath,amssymb,mathtools}
\usepackage{icomma}
\usepackage{xcolor}
\usepackage{colortbl}
\usepackage{tikz}
\usepackage{pgfplots}
\usepackage[most]{tcolorbox}
\usepackage{enumitem}
\usepackage{array}
\usepackage{booktabs}
\usepackage{multicol}
\usepackage{fancyhdr}
\usepackage{caption}
\captionsetup{labelformat=empty,justification=centering,font=small}
\usepackage[hidelinks]{hyperref}
\usetikzlibrary{arrows.meta,positioning,calc,angles,quotes,patterns,backgrounds,shapes.geometric}
\pgfplotsset{compat=1.18}
\usepackage[a4paper,margin=2.2cm,headheight=26pt,headsep=10pt]{geometry}

\definecolor{primary}{RGB}{219,54,77}
\definecolor{primarydark}{RGB}{150,28,46}
\definecolor{primarylight}{RGB}{250,224,228}
\definecolor{defcol}{RGB}{0,114,178}
\definecolor{thmcol}{RGB}{0,140,120}
\definecolor{excol}{RGB}{0,135,90}
\definecolor{methcol}{RGB}{90,90,100}
\definecolor{courbe}{RGB}{40,90,200}
\definecolor{courbeB}{RGB}{210,60,40}
\definecolor{courbeC}{RGB}{30,150,80}
\definecolor{axiscol}{RGB}{60,60,60}
\definecolor{gridcol}{RGB}{200,205,215}

\tcbset{boxbase/.style={enhanced,breakable,boxrule=0.9pt,arc=2.5pt,
   left=3mm,right=3mm,top=2.3mm,bottom=2.3mm,
   fonttitle=\bfseries,coltitle=white,toptitle=0.6mm,bottomtitle=0.6mm}}
\newtcolorbox{methode}[1][]{boxbase,colback=methcol!7,colframe=methcol,sharp corners,
  title={\textsc{Méthode}\ifx\relax#1\relax\relax\else\textmd{ --- \itshape #1}\fi}}
\newtcolorbox{exemple}[1][]{boxbase,colback=excol!5,colframe=excol,
  title={\textsc{Exemple}\ifx\relax#1\relax\relax\else\textmd{ : \itshape #1}\fi}}
\newtcolorbox{idee}[1][]{boxbase,colback=defcol!6,colframe=defcol,
  title={\textsc{L'essentiel}\ifx\relax#1\relax\relax\else\textmd{ : \itshape #1}\fi}}
\newtcolorbox{piege}[1][]{boxbase,colback=primary!7,colframe=primary,
  title={\textsc{Piège}\ifx\relax#1\relax\relax\else\textmd{ : \itshape #1}\fi}}

\pgfplotsset{repere/.style={axis lines=middle,
    axis line style={-{Stealth[length=2.2mm]},axiscol,line width=0.8pt},
    label style={font=\footnotesize},tick label style={font=\scriptsize},
    every axis x label/.style={at={(current axis.right of origin)},anchor=north west},
    every axis y label/.style={at={(current axis.above origin)},anchor=south east},
    xlabel={$x$},ylabel={$y$},grid=both,minor tick num=0,
    major grid style={gridcol,line width=0.4pt},samples=200,clip=false}}

\pagestyle{fancy}\fancyhf{}
\fancyhead[L]{\small\textcolor{primarydark}{\textbf{Fiche 4} — Les fonctions}}
\fancyhead[R]{\small\textcolor{primarydark}{\textsc{Carnet d'entraînement}}}
\fancyfoot[C]{\small\thepage}
\renewcommand{\headrulewidth}{0.8pt}
\renewcommand{\headrule}{\hbox to\headwidth{\color{primary}\leaders\hrule height \headrulewidth\hfill}}
\usepackage{titlesec}
\titleformat{\section}{\Large\bfseries\color{primarydark}}{\thesection}{0.6em}{}
\titleformat{\subsection}{\large\bfseries\color{primary}}{\thesubsection}{0.6em}{}

\newcommand{\R}{\mathbb{R}}\newcommand{\Z}{\mathbb{Z}}\newcommand{\N}{\mathbb{N}}
\newcommand{\Df}{D_f}
\newcommand{\croit}{\textcolor{thmcol}{\Large$\nearrow$}}
\newcommand{\decroit}{\textcolor{thmcol}{\Large$\searrow$}}

\newcommand{\exercice}[1]{\medskip\noindent{\color{primarydark}\bfseries\large Exercice #1}\par\nopagebreak\smallskip}
\newcommand{\titrecarnet}[3]{%
\begin{center}\begin{tikzpicture}
\node[rectangle,rounded corners=7pt,fill=primary,text=white,minimum width=15cm,
  minimum height=1.7cm,align=center,inner sep=8pt]
  {{\LARGE\bfseries #1}\\[2pt]{\normalsize #2}};
\end{tikzpicture}\end{center}
\begin{center}\itshape\color{primarydark}#3\end{center}\medskip}

% ---- RÈGLES D'USAGE DES MACROS ----
% * \croit et \decroit s'emploient en mode TEXTE (dans une cellule de tableau),
%   JAMAIS entre $...$ (ils contiennent déjà un $\nearrow$).
% * \titrecarnet{Titre}{Sous-titre}{Ligne en italique} : bandeau de titre.
% * \exercice{1} : titre d'exercice.
% * Boîtes : \begin{idee}...\end{idee}, \begin{exemple}[titre]...\end{exemple},
%   \begin{methode}[titre]...\end{methode}, \begin{piege}[titre]...\end{piege}.
% Le corps du document commence après \begin{document}.
\begin{document}
\titrecarnet{Thème 2 — Parité, symétries \& périodicité}{Tutoriel}{Repérer les symétries d'une courbe : axe, centre, parité, période.}

\section*{1. Parité : paire, impaire, ou ni l'un ni l'autre}

\begin{idee}[le test se fait toujours sur $f(-x)$]
On calcule $f(-x)$, puis on compare :
\begin{itemize}[leftmargin=1.5em,topsep=2pt]
  \item si $f(-x)=f(x)$ : $f$ est \textbf{paire} — courbe symétrique par rapport à l'\emph{axe des $y$} ;
  \item si $f(-x)=-f(x)$ : $f$ est \textbf{impaire} — courbe symétrique par rapport à l'\emph{origine} $(0;0)$ ;
  \item sinon : $f$ n'est \textbf{ni paire ni impaire} (le cas le plus fréquent !).
\end{itemize}
\end{idee}

\begin{piege}
« Paire » et « impaire » ne couvrent pas tous les cas. Pour prouver qu'une fonction n'est \emph{ni} l'une \emph{ni} l'autre, il suffit d'un contre-exemple : par exemple pour $f(x)=2x+1$, $f(1)=3$ mais $f(-1)=-1$, donc $f(-1)\neq f(1)$ et $f(-1)\neq -f(1)=-3$.
\end{piege}

\begin{center}
\begin{tikzpicture}
\begin{axis}[repere,width=6.6cm,height=4.4cm,xmin=-2.4,xmax=2.4,ymin=-0.5,ymax=4.2,
   xtick={-2,-1,1,2},ytick={1,2,3},samples=80,title={\footnotesize paire : $x^2+1$}]
  \addplot[courbe,line width=1.2pt,domain=-1.9:1.9]{x^2+1};
  \draw[dashed,gray](axis cs:0,-0.5)--(axis cs:0,4.2);
  \filldraw[courbeB](axis cs:1.3,2.69)circle(1.5pt);\filldraw[courbeB](axis cs:-1.3,2.69)circle(1.5pt);
\end{axis}
\end{tikzpicture}\hfill
\begin{tikzpicture}
\begin{axis}[repere,width=6.6cm,height=4.4cm,xmin=-2.1,xmax=2.1,ymin=-3.2,ymax=3.2,
   xtick={-2,-1,1,2},ytick={-2,2},samples=80,restrict y to domain=-3.2:3.2,title={\footnotesize impaire : $x^3$}]
  \addplot[courbe,line width=1.2pt,domain=-1.45:1.45]{x^3};
  \filldraw[courbeB](axis cs:1.2,1.728)circle(1.5pt);\filldraw[courbeB](axis cs:-1.2,-1.728)circle(1.5pt);
  \filldraw[black](axis cs:0,0)circle(1.4pt);
\end{axis}
\end{tikzpicture}
\end{center}

\section*{2. Axe et centre de symétrie}

\begin{idee}
\begin{itemize}[leftmargin=1.5em,topsep=2pt]
  \item \textbf{Axe de symétrie} $x=a$ : \ $f(a-x)=f(a+x)$ pour tout $x$. \ Pour une parabole $ax^2+bx+c$ : axe $x=-\dfrac{b}{2a}$.
  \item \textbf{Centre de symétrie} $(a;b)$ : \ $f(a-x)+f(a+x)=2b$ pour tout $x$ (invariance par demi-tour).
\end{itemize}
\end{idee}

\begin{exemple}
$f(x)=x^3+2$ admet le centre $(0;2)$ : $f(-x)+f(x)=(-x^3+2)+(x^3+2)=4=2\times 2$. \checkmark
\end{exemple}

\section*{3. Périodicité}

$f$ est \textbf{périodique de période $T$} si $f(x+T)=f(x)$ pour tout $x$ ($T$ = plus petit décalage $>0$ qui laisse la courbe inchangée). Le motif se répète tous les $T$.

\begin{idee}[périodes à connaître]
$\sin x$ et $\cos x$ : $T=2\pi$. \quad $\tan x$ : $T=\pi$. \quad $\sin(kx)$ et $\cos(kx)$ : $T=\dfrac{2\pi}{k}$. \quad $|\sin x|$ : $T=\pi$.
\end{idee}

\begin{exemple}
$\sin(2x)$ a pour période $T=\pi$ car $\sin\!\big(2(x+\pi)\big)=\sin(2x+2\pi)=\sin(2x)$.
\end{exemple}
\end{document}
